\documentclass[10pt,a4paper]{amsart}
\usepackage{amsmath,amssymb,amsthm,graphicx}
% Retained preamble for the Stacks Project chapter divisors.tex (Divisors).
% The chapter file pinned-original-source.tex carries only \input{preamble} on its first line, so no
% preamble text lives inside the chapter itself.  The Stacks Project uses one single shared preamble for
% every chapter; the snapshot stages it verbatim in sources/stacks-chow/preamble.tex.  The block below is
% copied mechanically, line for line, from that file's lines 14-18 (the Xy-pic setup) and lines 154-217
% (the native theorem declarations together with the mathematical macros).  Nothing was added, removed or
% redefined.  The wrapper class and the amsmath/amssymb/amsthm/graphicx base packages are supplied by the
% assembler, so the class/\externaldocument part of the shared preamble is not repeated here.
% For commutative diagrams we use Xy-pic
\usepackage[all]{xy}

% We use 2cell for 2-commutative diagrams.
\xyoption{2cell}
% Theorem environments.
%
\theoremstyle{plain}
\newtheorem{theorem}[subsection]{Theorem}
\newtheorem{proposition}[subsection]{Proposition}
\newtheorem{lemma}[subsection]{Lemma}

\theoremstyle{definition}
\newtheorem{definition}[subsection]{Definition}
\newtheorem{example}[subsection]{Example}
\newtheorem{exercise}[subsection]{Exercise}
\newtheorem{situation}[subsection]{Situation}

\theoremstyle{remark}
\newtheorem{remark}[subsection]{Remark}
\newtheorem{remarks}[subsection]{Remarks}

\numberwithin{equation}{subsection}

% Macros
%
\def\lim{\mathop{\mathrm{lim}}\nolimits}
\def\colim{\mathop{\mathrm{colim}}\nolimits}
\def\Spec{\mathop{\mathrm{Spec}}}
\def\Hom{\mathop{\mathrm{Hom}}\nolimits}
\def\Ext{\mathop{\mathrm{Ext}}\nolimits}
\def\SheafHom{\mathop{\mathcal{H}\!\mathit{om}}\nolimits}
\def\SheafExt{\mathop{\mathcal{E}\!\mathit{xt}}\nolimits}
\def\Sch{\mathit{Sch}}
\def\Mor{\mathop{\mathrm{Mor}}\nolimits}
\def\Ob{\mathop{\mathrm{Ob}}\nolimits}
\def\Sh{\mathop{\mathit{Sh}}\nolimits}
\def\NL{\mathop{N\!L}\nolimits}
\def\CH{\mathop{\mathrm{CH}}\nolimits}
\def\proetale{{pro\text{-}\acute{e}tale}}
\def\etale{{\acute{e}tale}}
\def\QCoh{\mathit{QCoh}}
\def\Ker{\mathop{\mathrm{Ker}}}
\def\Im{\mathop{\mathrm{Im}}}
\def\Coker{\mathop{\mathrm{Coker}}}
\def\Coim{\mathop{\mathrm{Coim}}}

% Boxtimes
%
\DeclareMathSymbol{\boxtimes}{\mathbin}{AMSa}{"02}

%
% Macros for moduli stacks/spaces
%
\def\QCohstack{\mathcal{QC}\!\mathit{oh}}
\def\Cohstack{\mathcal{C}\!\mathit{oh}}
\def\Spacesstack{\mathcal{S}\!\mathit{paces}}
\def\Quotfunctor{\mathrm{Quot}}
\def\Hilbfunctor{\mathrm{Hilb}}
\def\Curvesstack{\mathcal{C}\!\mathit{urves}}
\def\Polarizedstack{\mathcal{P}\!\mathit{olarized}}
\def\Complexesstack{\mathcal{C}\!\mathit{omplexes}}
% \Pic is the operator that assigns to X its picard group, usage \Pic(X)
% \Picardstack_{X/B} denotes the Picard stack of X over B
% \Picardfunctor_{X/B} denotes the Picard functor of X over B
\def\Pic{\mathop{\mathrm{Pic}}\nolimits}
\def\Picardstack{\mathcal{P}\!\mathit{ic}}
\def\Picardfunctor{\mathrm{Pic}}
\def\Deformationcategory{\mathcal{D}\!\mathit{ef}}

% --- appended 2026-10-12 (Chief): the project own definitions of the reference/slogan/history
% display environments, copied mechanically from the shared Stacks preamble lines 8-12
% (\usepackage{verbatim} plus the three \newenvironment lines, which the chapter preamble keeps
% verbatim from the project preamble). Nothing added, removed or redefined: these environments
% carry no mathematical content and the project itself typesets them as comments. Without these
% three lines a counted statement whose own text uses them fails with Environment slogan undefined.
\usepackage{verbatim}
\newenvironment{reference}{\comment}{\endcomment}
\newenvironment{slogan}{\comment}{\endcomment}
\newenvironment{history}{\comment}{\endcomment}
\usepackage{cleveref}
\crefname{originallemma}{original lemma}{original lemmas}
\Crefname{originallemma}{Original lemma}{Original lemmas}
\makeatletter
\newlabel{algebra-lemma-dimension-at-a-point-finite-type-field}{{00P1}{source}{}{sourcepointer-1}{}}
\newlabel{algebra-lemma-dimension-at-a-point-finite-type-field@cref}{{[originallemma][0][]00P1}{[1][0][]source}{}{}{}}
\makeatother
\crefname{originallemma}{original lemma}{original lemmas}
\Crefname{originallemma}{Original lemma}{Original lemmas}
\makeatletter
\newlabel{algebra-lemma-local-isomorphism}{{00QR}{source}{}{sourcepointer-2}{}}
\newlabel{algebra-lemma-local-isomorphism@cref}{{[originallemma][0][]00QR}{[1][0][]source}{}{}{}}
\makeatother
\crefname{originallemma}{original lemma}{original lemmas}
\Crefname{originallemma}{Original lemma}{Original lemmas}
\makeatletter
\newlabel{algebra-lemma-mod-injective-general}{{046Y}{source}{}{sourcepointer-3}{}}
\newlabel{algebra-lemma-mod-injective-general@cref}{{[originallemma][0][]046Y}{[1][0][]source}{}{}{}}
\makeatother
\crefname{originallemma}{original lemma}{original lemmas}
\Crefname{originallemma}{Original lemma}{Original lemmas}
\makeatletter
\newlabel{algebra-lemma-one-equation}{{00KW}{source}{}{sourcepointer-4}{}}
\newlabel{algebra-lemma-one-equation@cref}{{[originallemma][0][]00KW}{[1][0][]source}{}{}{}}
\makeatother
\crefname{originallemma}{original lemma}{original lemmas}
\Crefname{originallemma}{Original lemma}{Original lemmas}
\makeatletter
\newlabel{algebra-lemma-quasi-regular-regular}{{061S}{source}{}{sourcepointer-5}{}}
\newlabel{algebra-lemma-quasi-regular-regular@cref}{{[originallemma][0][]061S}{[1][0][]source}{}{}{}}
\makeatother
\crefname{originallemma}{original lemma}{original lemmas}
\Crefname{originallemma}{Original lemma}{Original lemmas}
\makeatletter
\newlabel{algebra-lemma-regular-sequence-in-neighbourhood}{{061L}{source}{}{sourcepointer-6}{}}
\newlabel{algebra-lemma-regular-sequence-in-neighbourhood@cref}{{[originallemma][0][]061L}{[1][0][]source}{}{}{}}
\makeatother
\crefname{originallemma}{original lemma}{original lemmas}
\Crefname{originallemma}{Original lemma}{Original lemmas}
\makeatletter
\newlabel{lemma-fibre-Cartier}{{062Y}{source}{}{sourcepointer-7}{}}
\newlabel{lemma-fibre-Cartier@cref}{{[originallemma][0][]062Y}{[1][0][]source}{}{}{}}
\makeatother
\crefname{originallemma}{original lemma}{original lemmas}
\Crefname{originallemma}{Original lemma}{Original lemmas}
\makeatletter
\newlabel{lemma-relative-regular-immersion-flat-in-neighbourhood}{{063U}{source}{}{sourcepointer-8}{}}
\newlabel{lemma-relative-regular-immersion-flat-in-neighbourhood@cref}{{[originallemma][0][]063U}{[1][0][]source}{}{}{}}
\makeatother
\crefname{originallemma}{original lemma}{original lemmas}
\Crefname{originallemma}{Original lemma}{Original lemmas}
\makeatletter
\newlabel{morphisms-lemma-composition-finite-presentation}{{01TR}{source}{}{sourcepointer-9}{}}
\newlabel{morphisms-lemma-composition-finite-presentation@cref}{{[originallemma][0][]01TR}{[1][0][]source}{}{}{}}
\makeatother
\crefname{originallemma}{original lemma}{original lemmas}
\Crefname{originallemma}{Original lemma}{Original lemmas}
\makeatletter
\newlabel{morphisms-lemma-finite-presentation-permanence}{{02FV}{source}{}{sourcepointer-10}{}}
\newlabel{morphisms-lemma-finite-presentation-permanence@cref}{{[originallemma][0][]02FV}{[1][0][]source}{}{}{}}
\makeatother
\title{A closed subscheme whose fibres are quasi-regular immersions is locally contained in a relative quasi-regular immersion, and agrees with one under flatness or under the codimension count}
\author{The Stacks Project Authors}
\date{stacks-project:divisors.tex:local-snapshot-20261012; GNU FDL 1.2 or later}
\begin{document}
\maketitle
\noindent\textit{Editorial note.} All delivered mathematical text is the author's own, copied from the pinned Stacks Project chapter file pinned-original-source.tex; no statement or proof step was rewritten, reordered or removed.\par\medskip
\noindent\textit{Editorial note.} Following the multi-claim convention for this source, source\_url is the source-level URL https://stacks.math.columbia.edu/ and the source label lemma-fibre-quasi-regular is carried in claim\_key and here rather than in a per-claim tag URL.\par\medskip
\noindent\textit{Editorial note.} The preamble is the chapter's shared Stacks Project preamble, retained as sources/stacks-divisors/retained-preamble.tex: a line-for-line excerpt of the shared Stacks preamble staged with this snapshot. No macro was invented or redefined, and this card adds no preamble text of its own.\par\medskip
\noindent\textit{Editorial note.} Statement span 5491-5515 and proof span 5517-5594 come from the halfA pool card operations/pool-cards/stacks-divisors-halfA.json and were re-read in the pinned chapter file before this card was written.\par\medskip
\noindent\textit{Editorial note.} The proof of part (2) is written with the letters of the preceding lemma (it speaks of a closed immersion between the two constructed subschemes and of the corresponding local rings); the wording is the author's and is delivered unchanged, and no letter was normalised by us.\par\medskip
\noindent\textit{Editorial note.} The proof of part (3) contains the author's parenthesis '(some details omitted)' after the statement that the change-of-basis matrix is invertible modulo the maximal ideal; that phrase is kept verbatim.\par\medskip
\noindent\textit{Editorial note.} 10 cross-references occur inside the counted spans; they are listed in external\_references as navigation pointers only.\par\medskip
\noindent\textit{Allowed standard prerequisites (editorial).}\begin{itemize}
\item The relation between quasi-regularity and regularity of a sequence after localization, and the persistence of a regular sequence in a neighbourhood of a prime. Cited at proof lines 5533-5536. Turns the quasi-regular sequence generating the fibre ideal into a regular sequence on a neighbourhood, which is what the iterative construction needs.
\item The fibrewise criterion for a closed subscheme to be a relative effective Cartier divisor, applied one equation at a time, cited from the same chapter. Cited at proof lines 5537-5543. Produces each successive member of the filtration by vanishing loci of the chosen sequence.
\item A surjective map of flat local algebras essentially of finite presentation that becomes an isomorphism after dividing by the maximal ideal of the base is an isomorphism, and the local criterion for an isomorphism in a neighbourhood. Cited at proof lines 5554-5559. Upgrades the containment of part (1) to equality in the flat case of part (2).
\item The dimension formula at a point of a finite type scheme over a field and the drop of dimension by one on cutting out a nonzerodivisor. Cited at proof lines 5567-5573. Shows that the number of equations cutting out the fibre in the hypothesis of part (3) equals the length of the regular sequence constructed in part (1).
\end{itemize}\medskip
\ifdefined\hypertarget\hypertarget{sourcepointer-1}{}\fi
\noindent\textit{Source reference.} Original lemma 00P1: focused-seeds/AGP/sources/stacks-snapshot-20261011/tags.txt, source lines 915--915. Same-chapter reference used in the retained author text at source line(s) 5576. It is supplied only as a navigation pointer so that the retained text compiles; the counted item delivers the author's own statement and the author's complete proof, and no part of the referenced item is copied in or substituted for the counted argument.\par\medskip
\ifdefined\hypertarget\hypertarget{sourcepointer-2}{}\fi
\noindent\textit{Source reference.} Original lemma 00QR: focused-seeds/AGP/sources/stacks-snapshot-20261011/tags.txt, source lines 977--977. Same-chapter reference used in the retained author text at source line(s) 4274, 5560. It is supplied only as a navigation pointer so that the retained text compiles; the counted item delivers the author's own statement and the author's complete proof, and no part of the referenced item is copied in or substituted for the counted argument.\par\medskip
\ifdefined\hypertarget\hypertarget{sourcepointer-3}{}\fi
\noindent\textit{Source reference.} Original lemma 046Y: focused-seeds/AGP/sources/stacks-snapshot-20261011/tags.txt, source lines 5489--5489. Same-chapter reference used in the retained author text at source line(s) 4271, 5557. It is supplied only as a navigation pointer so that the retained text compiles; the counted item delivers the author's own statement and the author's complete proof, and no part of the referenced item is copied in or substituted for the counted argument.\par\medskip
\ifdefined\hypertarget\hypertarget{sourcepointer-4}{}\fi
\noindent\textit{Source reference.} Original lemma 00KW: focused-seeds/AGP/sources/stacks-snapshot-20261011/tags.txt, source lines 766--766. Same-chapter reference used in the retained author text at source line(s) 3036, 3378, 5578. It is supplied only as a navigation pointer so that the retained text compiles; the counted item delivers the author's own statement and the author's complete proof, and no part of the referenced item is copied in or substituted for the counted argument.\par\medskip
\ifdefined\hypertarget\hypertarget{sourcepointer-5}{}\fi
\noindent\textit{Source reference.} Original lemma 061S: focused-seeds/AGP/sources/stacks-snapshot-20261011/tags.txt, source lines 7771--7771. Same-chapter reference used in the retained author text at source line(s) 4873, 5318, 5358, 5533. It is supplied only as a navigation pointer so that the retained text compiles; the counted item delivers the author's own statement and the author's complete proof, and no part of the referenced item is copied in or substituted for the counted argument.\par\medskip
\ifdefined\hypertarget\hypertarget{sourcepointer-6}{}\fi
\noindent\textit{Source reference.} Original lemma 061L: focused-seeds/AGP/sources/stacks-snapshot-20261011/tags.txt, source lines 7765--7765. Same-chapter reference used in the retained author text at source line(s) 2970, 4250, 4869, 5359, 5534, 5635. It is supplied only as a navigation pointer so that the retained text compiles; the counted item delivers the author's own statement and the author's complete proof, and no part of the referenced item is copied in or substituted for the counted argument.\par\medskip
\ifdefined\hypertarget\hypertarget{sourcepointer-7}{}\fi
\noindent\textit{Source reference.} Original lemma 062Y: focused-seeds/AGP/sources/stacks-snapshot-20261011/tags.txt, source lines 7812--7812. Same-chapter reference used in the retained author text at source line(s) 5361, 5538. It is supplied only as a navigation pointer so that the retained text compiles; the counted item delivers the author's own statement and the author's complete proof, and no part of the referenced item is copied in or substituted for the counted argument.\par\medskip
\ifdefined\hypertarget\hypertarget{sourcepointer-8}{}\fi
\noindent\textit{Source reference.} Original lemma 063U: focused-seeds/AGP/sources/stacks-snapshot-20261011/tags.txt, source lines 7843--7843. Same-chapter reference used in the retained author text at source line(s) 5301, 5387, 5483, 5511. It is supplied only as a navigation pointer so that the retained text compiles; the counted item delivers the author's own statement and the author's complete proof, and no part of the referenced item is copied in or substituted for the counted argument.\par\medskip
\ifdefined\hypertarget\hypertarget{sourcepointer-9}{}\fi
\noindent\textit{Source reference.} Original lemma 01TR: focused-seeds/AGP/sources/stacks-snapshot-20261011/tags.txt, source lines 2384--2384. Same-chapter reference used in the retained author text at source line(s) 4289, 5592. It is supplied only as a navigation pointer so that the retained text compiles; the counted item delivers the author's own statement and the author's complete proof, and no part of the referenced item is copied in or substituted for the counted argument.\par\medskip
\ifdefined\hypertarget\hypertarget{sourcepointer-10}{}\fi
\noindent\textit{Source reference.} Original lemma 02FV: focused-seeds/AGP/sources/stacks-snapshot-20261011/tags.txt, source lines 3206--3206. Same-chapter reference used in the retained author text at source line(s) 4290, 5593. It is supplied only as a navigation pointer so that the retained text compiles; the counted item delivers the author's own statement and the author's complete proof, and no part of the referenced item is copied in or substituted for the counted argument.\par\medskip
\begin{lemma}
\label{lemma-fibre-quasi-regular}
Let $\varphi : X \to S$ be a flat morphism which is locally of finite
presentation. Let $T \subset X$ be a closed subscheme.
Let $x \in T$ with image $s \in S$.
\begin{enumerate}
\item If $T_s \subset X_s$ is a quasi-regular immersion
in a neighbourhood of $x$, then there exists an open
$U \subset X$ and a relative quasi-regular immersion
$Z \subset U$ such that $Z_s = T_s \cap U_s$ and $T \cap U \subset Z$.
\item If $T_s \subset X_s$ is a quasi-regular immersion
in a neighbourhood of $x$, the morphism $T \to X$ is of finite
presentation, and $T \to S$ is flat at $x$, then we can choose $U$ and
$Z$ as in (1) such that $T \cap U = Z$.
\item If $T_s \subset X_s$ is a quasi-regular immersion in a neighbourhood
of $x$, and $T$ is cut out by $c$ equations in a neighbourhood of $x$,
where $c = \dim_x(X_s) - \dim_x(T_s)$, then we can choose $U$ and $Z$ as in (1)
such that $T \cap U = Z$.
\end{enumerate}
In each case $Z \to U$ is a regular immersion by
Lemma \ref{lemma-relative-regular-immersion-flat-in-neighbourhood}.
In particular, if $T \to S$ is locally of finite presentation and flat and
all fibres $T_s \subset X_s$ are quasi-regular immersions, then
$T \to X$ is a relative quasi-regular immersion.
\end{lemma}

\begin{proof}
Choose affine open neighbourhoods $\Spec(A)$ of $s$ and
$\Spec(B)$ of $x$ such that
$\varphi(\Spec(B)) \subset \Spec(A)$.
Let $\mathfrak p \subset A$ be the prime ideal corresponding to $s$.
Let $\mathfrak q \subset B$ be the prime ideal corresponding to $x$.
Let $I \subset B$ be the ideal corresponding to $T$.
By the initial assumption of the lemma we know that
$A \to B$ is flat and of finite presentation.
The assumption in (1) means that, after shrinking $\Spec(B)$, we may
assume $I(B \otimes_A \kappa(\mathfrak p))$ is generated by a
quasi-regular sequence of elements. After possibly localizing $B$
at some $g \in B$, $g \not \in \mathfrak q$ we may assume there
exist $f_1, \ldots, f_r \in I$ which map to a quasi-regular
sequence in $B \otimes_A \kappa(\mathfrak p)$ which generates
$I(B \otimes_A \kappa(\mathfrak p))$. By
Algebra, Lemmas \ref{algebra-lemma-quasi-regular-regular} and
\ref{algebra-lemma-regular-sequence-in-neighbourhood}
we may assume after another localization that
$f_1, \ldots, f_r \in I$ form a regular
sequence in $B \otimes_A \kappa(\mathfrak p)$. By
Lemma \ref{lemma-fibre-Cartier}
it follows that $Z_1 = V(f_1) \subset \Spec(B)$
is a relative effective Cartier divisor, again after possibly
localizing $B$. Applying the same lemma again, but now to
$Z_2 = V(f_1, f_2) \subset Z_1$ we see that $Z_2 \subset Z_1$
is a relative effective Cartier divisor. And so on until one
reaches $Z = Z_n = V(f_1, \ldots, f_n)$. Then
$Z \to \Spec(B)$ is a regular immersion and $Z$ is
flat over $S$, in particular $Z \to \Spec(B)$ is
a relative quasi-regular immersion over $\Spec(A)$.
This proves (1).

\medskip\noindent
To see (2) consider the closed immersion $Z \to D$. The surjective
ring map $u : \mathcal{O}_{D, x} \to \mathcal{O}_{Z, x}$
is a map of flat local $\mathcal{O}_{S, s}$-algebras which
are essentially of finite presentation, and which becomes an
isomorphisms after dividing by $\mathfrak m_s$. Hence it is
an isomorphism, see
Algebra, Lemma \ref{algebra-lemma-mod-injective-general}.
It follows that $Z \to D$ is an isomorphism in a neighbourhood
of $x$, see
Algebra, Lemma \ref{algebra-lemma-local-isomorphism}.

\medskip\noindent
To see (3), after possibly shrinking $U$ we may assume that
the ideal of $Z$ is generated by a regular sequence $f_1, \ldots, f_r$
(see our construction of $Z$ above)
and the ideal of $T$ is generated by $g_1, \ldots, g_c$.
We claim that $c = r$. Namely,
\begin{align*}
\dim_x(X_s) & = \dim(\mathcal{O}_{X_s, x}) +
\text{trdeg}_{\kappa(s)}(\kappa(x)), \\
\dim_x(T_s) & = \dim(\mathcal{O}_{T_s, x}) +
\text{trdeg}_{\kappa(s)}(\kappa(x)), \\
\dim(\mathcal{O}_{X_s, x}) & = \dim(\mathcal{O}_{T_s, x}) + r
\end{align*}
the first two equalities by
Algebra, Lemma \ref{algebra-lemma-dimension-at-a-point-finite-type-field}
and the second by $r$ times applying
Algebra, Lemma \ref{algebra-lemma-one-equation}.
As $T \subset Z$ we see that
$f_i = \sum b_{ij} g_j$. But the ideals of $Z$ and $T$ cut out the same
quasi-regular closed subscheme of $X_s$ in a neighbourhood of $x$. Hence
the matrix $(b_{ij}) \bmod \mathfrak m_x$ is invertible (some details
omitted). Hence $(b_{ij})$ is invertible in an
open neighbourhood of $x$. In other words,
$T \cap U = Z$ after shrinking $U$.

\medskip\noindent
The final statements of the lemma follow immediately from
part (2), combined with the fact that $Z \to S$
is locally of finite presentation if and only if $Z \to X$ is
of finite presentation, see
Morphisms, Lemmas \ref{morphisms-lemma-composition-finite-presentation} and
\ref{morphisms-lemma-finite-presentation-permanence}.
\end{proof}

\end{document}
